\documentclass[11pt,a4paper,landscape]{article}
\usepackage[T1]{fontenc}
\usepackage[utf8]{inputenc}
\usepackage[provide=*,lithuanian]{babel}
\usepackage{lmodern}
\usepackage{amsmath,amssymb}
\usepackage[margin=0pt]{geometry}
\usepackage{tikz}
\pagestyle{empty}
\setlength{\parindent}{0pt}
\setlength{\parskip}{0pt}

% Kompiliavimas: pdfLaTeX.
% Spausdinti horizontaliame A4 lape, mastelis 100 % (tikrasis dydis).
% Perkirpus per vidurį gaunami du statūs A5 dydžio lapeliai.
% A4 pusės plotis yra 148,5 mm; standartinio A5 nominalus plotis – 148 mm.

% Tuščias stačiakampis: plotis ir aukštis.
\newcommand{\remelis}[2]{%
  \tikz[baseline={([yshift=-0.6ex]current bounding box.center)}]\draw[line width=0.5pt]
    (0,0) rectangle (#1,#2);}

% Viena užduotis. Numeris ir sąlyga įrašomi kaip du argumentai.
% Jei reikia daugiau vietos sąlygai, didinkite 24 mm eilutės aukštį.
% Keičiant užduočių skaičių patikrinkite, ar turinys telpa į lapelį.
\newcommand{\uzduotis}[2]{%
  \begin{minipage}[t][24mm][t]{\linewidth}
    \textbf{#1.}\enspace #2\par
    \vspace{3mm}
    \remelis{42mm}{10mm}\par
  \end{minipage}\par
}

% Vieno A5 lapelio turinys.
% #1 – varianto pavadinimas; #2 – užduočių sąrašas.
\newcommand{\lapelis}[2]{%
  \begin{minipage}[t][190mm][t]{128.5mm}
    \small
    \textbf{Mokinio kodas}\enspace\remelis{32mm}{8mm}%
    \hfill\textbf{#1}\par
    \vspace{6mm}
    {\centering\normalsize\textbf{Trumpa apklausa}\par}
    \vspace{2mm}
    {\centering\small Į rėmelius įrašykite tik atsakymus.\par}
    \vspace{7mm}
    \normalsize
    #2
  \end{minipage}%
}

\begin{document}
\null
\begin{tikzpicture}[remember picture,overlay]
  % Punktyrinė kirpimo linija tiksliai A4 lapo viduryje.
  \draw[dash pattern=on 2pt off 3pt,line width=0.4pt]
    (current page.north) -- (current page.south);

  % Kairysis A5 lapelis: 10 mm paraštės.
  \node[anchor=north west,inner sep=0pt]
    at ([xshift=10mm,yshift=-10mm]current page.north west) {%
      \lapelis{I variantas}{%
        \begin{minipage}[t]{60mm}
          \uzduotis{1}{Užduoties sąlyga.}
          \uzduotis{2}{Užduoties sąlyga.}
          \uzduotis{3}{Užduoties sąlyga.}
          \uzduotis{4}{Užduoties sąlyga.}
          \uzduotis{5}{Užduoties sąlyga.}
          \uzduotis{6}{Užduoties sąlyga.}
        \end{minipage}\hfill
        \begin{minipage}[t]{60mm}
          \uzduotis{7}{Užduoties sąlyga.}
          \uzduotis{8}{Užduoties sąlyga.}
          \uzduotis{9}{Užduoties sąlyga.}
          \uzduotis{10}{Užduoties sąlyga.}
          \uzduotis{11}{Užduoties sąlyga.}
          \uzduotis{12}{Užduoties sąlyga.}
        \end{minipage}
      }%
    };

  % Dešinysis A5 lapelis.
  \node[anchor=north west,inner sep=0pt]
    at ([xshift=10mm,yshift=-10mm]current page.north) {%
      \lapelis{II variantas}{%
        \begin{minipage}[t]{60mm}
          \uzduotis{1}{Užduoties sąlyga.}
          \uzduotis{2}{Užduoties sąlyga.}
          \uzduotis{3}{Užduoties sąlyga.}
          \uzduotis{4}{Užduoties sąlyga.}
          \uzduotis{5}{Užduoties sąlyga.}
          \uzduotis{6}{Užduoties sąlyga.}
        \end{minipage}\hfill
        \begin{minipage}[t]{60mm}
          \uzduotis{7}{Užduoties sąlyga.}
          \uzduotis{8}{Užduoties sąlyga.}
          \uzduotis{9}{Užduoties sąlyga.}
          \uzduotis{10}{Užduoties sąlyga.}
          \uzduotis{11}{Užduoties sąlyga.}
          \uzduotis{12}{Užduoties sąlyga.}
        \end{minipage}
      }%
    };
\end{tikzpicture}
\end{document}
